
\documentclass{article}
\usepackage{colortbl}
\usepackage{makecell}
\usepackage{multirow}
\usepackage{supertabular}

\begin{document}

\newcounter{utterance}

\centering \large Interaction Transcript for game `dond', experiment `coop', episode 2 with nemotron{-}3{-}ultra.
\vspace{24pt}

{ \footnotesize  \setcounter{utterance}{1}
\setlength{\tabcolsep}{0pt}
\begin{supertabular}{c@{$\;$}|p{.15\linewidth}@{}p{.15\linewidth}p{.15\linewidth}p{.15\linewidth}p{.15\linewidth}p{.15\linewidth}}
    \# & $\;$A & \multicolumn{4}{c}{Game Master} & $\;\:$B\\
    \hline

    \theutterance \stepcounter{utterance}  
    & & \multicolumn{4}{p{0.6\linewidth}}{
        \cellcolor[rgb]{0.9,0.9,0.9}{
            \makecell[{{p{\linewidth}}}]{
                \texttt{\tiny{[P1$\langle$GM]}}
                \texttt{Du spiller et kooperativt forhandlingsspil, hvor du skal enes med en anden spiller om, hvordan du deler et sæt genstande mellem jer.} \\
\\ 
\texttt{Regler:} \\
\texttt{(a) Du og den anden spiller får et sæt genstande. Du får også en hemlig værdifunktion, der viser, hvor meget du værdisætter hver type objekt.} \\
\texttt{(b) Du bytter beskeder med den anden spiller for at blive enig om, hvem der får hvad. Du kan sende maksimalt 5 beskeder hver eller afslutte spillet tidligt ved at fremsætte et hemmeligt forslag når som helst.} \\
\texttt{(c) Du skal indsende et hemmeligt forslag, der angiver de genstande, du vil have, formateret i firkantede parenteser som følger: "{[}Forslag: <antal> <objektnavn>, <antal> <objektnavn>, <...>{]}"} \\
\texttt{(d) Hvis dine forslag er komplementære, dvs. der er nok genstande til at opfylde begge forslag, tildeles hver spiller en score baseret på summen af værdierne for de genstande, hun modtog. Ellers får du begge nul point.} \\
\texttt{(e) Dit mål er at maksimere summen af din score og den anden spillers score.} \\
\\ 
\texttt{Lad os starte.} \\
\\ 
\texttt{Sættet af tilgængelige genstande er:} \\
\\ 
\texttt{1 pose, 3 lamper, 1 blomst} \\
\\ 
\texttt{Din hemmelige værdifunktion er:} \\
\\ 
\texttt{pose: 3, lampe: 0, blomst: 7} \\
\\ 
\texttt{VIGTIG: Dine beskeder, medmindre det er et hemmeligt forslag, bliver sendt direkte til den anden spiller, så inkluder ikke noget svar på reglerne eller tekst, der annoncerer din besked. Brug det angivne format for at afgive et hemmeligt forslag. Brug ikke firkantede parenteser, når du kommunikerer med den anden spiller, eller det vil blive fortolket som dit hemmelige forslag.} \\
            }
        }
    }
    & \\ \\

    \theutterance \stepcounter{utterance}  
    & \multicolumn{4}{p{0.6\linewidth}}{
        \cellcolor[rgb]{0.9,0.9,0.9}{
            \makecell[{{p{\linewidth}}}]{
                \texttt{\tiny{[P1$\rangle$GM]}}
                \texttt{Hej! Jeg værdisætter blomsten meget højt (7 point), mens lampérne ikke er værd nogen ting for mig (0 point). Posen er lidt værd (3 point).} \\
\\ 
\texttt{Hvad værdisætter du mest? Hvis du gerne vil have lampérne og/eller posen, kan vi nemt finde en aftale, hvor vi begge får det, vi ønsker.} \\
            }
        }
    }
    & & \\ \\

    \theutterance \stepcounter{utterance}  
    & & & \multicolumn{2}{p{0.3\linewidth}}{
        \cellcolor[rgb]{0.9,0.9,0.9}{
            \makecell[{{p{\linewidth}}}]{
                \texttt{\tiny{[GM$|$GM]}}
                \texttt{continue} \\
            }
        }
    }
    & & \\ \\

    \theutterance \stepcounter{utterance}  
    & & & \multicolumn{2}{p{0.3\linewidth}}{
        \cellcolor[rgb]{0.9,0.9,0.9}{
            \makecell[{{p{\linewidth}}}]{
                \texttt{\tiny{[GM$|$GM]}}
                \texttt{Hej! Jeg værdisætter blomsten meget højt (7 point), mens lampérne ikke er værd nogen ting for mig (0 point). Posen er lidt værd (3 point).} \\
\\ 
\texttt{Hvad værdisætter du mest? Hvis du gerne vil have lampérne og/eller posen, kan vi nemt finde en aftale, hvor vi begge får det, vi ønsker.} \\
            }
        }
    }
    & & \\ \\

    \theutterance \stepcounter{utterance}  
    & & \multicolumn{4}{p{0.6\linewidth}}{
        \cellcolor[rgb]{0.9,0.9,0.9}{
            \makecell[{{p{\linewidth}}}]{
                \texttt{\tiny{[GM$\rangle$P2]}}
                \texttt{Du spiller et kooperativt forhandlingsspil, hvor du skal enes med en anden spiller om, hvordan du deler et sæt genstande mellem jer.} \\
\\ 
\texttt{Regler:} \\
\texttt{(a) Du og den anden spiller får et sæt genstande. Du får også en hemlig værdifunktion, der viser, hvor meget du værdisætter hver type objekt.} \\
\texttt{(b) Du bytter beskeder med den anden spiller for at blive enig om, hvem der får hvad. Du kan sende maksimalt 5 beskeder hver eller afslutte spillet tidligt ved at fremsætte et hemmeligt forslag når som helst.} \\
\texttt{(c) Du skal indsende et hemmeligt forslag, der angiver de genstande, du vil have, formateret i firkantede parenteser som følger: "{[}Forslag: <antal> <objektnavn>, <antal> <objektnavn>, <...>{]}"} \\
\texttt{(d) Hvis dine forslag er komplementære, dvs. der er nok genstande til at opfylde begge forslag, tildeles hver spiller en score baseret på summen af værdierne for de genstande, hun modtog. Ellers får du begge nul point.} \\
\texttt{(e) Dit mål er at maksimere summen af din score og den anden spillers score.} \\
\\ 
\texttt{Lad os starte.} \\
\\ 
\texttt{Sættet af tilgængelige genstande er:} \\
\\ 
\texttt{1 pose, 3 lamper, 1 blomst} \\
\\ 
\texttt{Din hemmelige værdifunktion er:} \\
\\ 
\texttt{pose: 1, lampe: 2, blomst: 3} \\
\\ 
\texttt{VIGTIG: Dine beskeder, medmindre det er et hemmeligt forslag, bliver sendt direkte til den anden spiller, så inkluder ikke noget svar på reglerne eller tekst, der annoncerer din besked. Brug det angivne format for at afgive et hemmeligt forslag. Brug ikke firkantede parenteser, når du kommunikerer med den anden spiller, eller det vil blive fortolket som dit hemmelige forslag.} \\
\\ 
\texttt{Hej! Jeg værdisætter blomsten meget højt (7 point), mens lampérne ikke er værd nogen ting for mig (0 point). Posen er lidt værd (3 point).} \\
\\ 
\texttt{Hvad værdisætter du mest? Hvis du gerne vil have lampérne og/eller posen, kan vi nemt finde en aftale, hvor vi begge får det, vi ønsker.} \\
            }
        }
    }
    & \\ \\

    \theutterance \stepcounter{utterance}  
    & & & \multicolumn{4}{p{0.6\linewidth}}{
        \cellcolor[rgb]{0.9,0.9,0.9}{
            \makecell[{{p{\linewidth}}}]{
                \texttt{\tiny{[GM$\langle$P2]}}
                \texttt{Jeg værdisætter blomsten højest (3 point), derefter lampérne (2 point hver) og posen mindst (1 point).} \\
\\ 
\texttt{Da du siger, at du værdisætter blomsten mest (og at lampérne er 0 point for dig), er der god grundlag for en aftale: Du tager blomsten og måske posen, jeg tager de tre lampér. Sådan får vi begge høje point.} \\
\\ 
\texttt{Hvad tænker du om, at du får blomsten (+ posen hvis du vil) og jeg får de tre lampér?} \\
            }
        }
    }
     \\ \\

    \theutterance \stepcounter{utterance}  
    & & & \multicolumn{2}{p{0.3\linewidth}}{
        \cellcolor[rgb]{0.9,0.9,0.9}{
            \makecell[{{p{\linewidth}}}]{
                \texttt{\tiny{[GM$|$GM]}}
                \texttt{continue} \\
            }
        }
    }
    & & \\ \\

    \theutterance \stepcounter{utterance}  
    & & & \multicolumn{2}{p{0.3\linewidth}}{
        \cellcolor[rgb]{0.9,0.9,0.9}{
            \makecell[{{p{\linewidth}}}]{
                \texttt{\tiny{[GM$|$GM]}}
                \texttt{Jeg værdisætter blomsten højest (3 point), derefter lampérne (2 point hver) og posen mindst (1 point).} \\
\\ 
\texttt{Da du siger, at du værdisætter blomsten mest (og at lampérne er 0 point for dig), er der god grundlag for en aftale: Du tager blomsten og måske posen, jeg tager de tre lampér. Sådan får vi begge høje point.} \\
\\ 
\texttt{Hvad tænker du om, at du får blomsten (+ posen hvis du vil) og jeg får de tre lampér?} \\
            }
        }
    }
    & & \\ \\

    \theutterance \stepcounter{utterance}  
    & & \multicolumn{4}{p{0.6\linewidth}}{
        \cellcolor[rgb]{0.9,0.9,0.9}{
            \makecell[{{p{\linewidth}}}]{
                \texttt{\tiny{[P1$\langle$GM]}}
                \texttt{Jeg værdisætter blomsten højest (3 point), derefter lampérne (2 point hver) og posen mindst (1 point).} \\
\\ 
\texttt{Da du siger, at du værdisætter blomsten mest (og at lampérne er 0 point for dig), er der god grundlag for en aftale: Du tager blomsten og måske posen, jeg tager de tre lampér. Sådan får vi begge høje point.} \\
\\ 
\texttt{Hvad tænker du om, at du får blomsten (+ posen hvis du vil) og jeg får de tre lampér?} \\
            }
        }
    }
    & \\ \\

    \theutterance \stepcounter{utterance}  
    & \multicolumn{4}{p{0.6\linewidth}}{
        \cellcolor[rgb]{0.9,0.9,0.9}{
            \makecell[{{p{\linewidth}}}]{
                \texttt{\tiny{[P1$\rangle$GM]}}
                \texttt{Det lyder som en god plan. Jeg er meget glad for at tage blomsten, da den giver mig 7 point. Posen giver mig 3 point, så jeg tager den med gerne, men det er ikke nødvendigt for at nå en god samlet score.} \\
\\ 
\texttt{Vi enes om: **Jeg får blomsten (og posen), du får de tre lampér.**} \\
\\ 
\texttt{Jeg sender nu mit hemmelige forslag.} \\
\\ 
\texttt{{[}Forslag: 1 blomst, 1 pose{]}} \\
            }
        }
    }
    & & \\ \\

    \theutterance \stepcounter{utterance}  
    & & & \multicolumn{2}{p{0.3\linewidth}}{
        \cellcolor[rgb]{0.9,0.9,0.9}{
            \makecell[{{p{\linewidth}}}]{
                \texttt{\tiny{[GM$|$GM]}}
                \texttt{proposal} \\
            }
        }
    }
    & & \\ \\

    \theutterance \stepcounter{utterance}  
    & & & \multicolumn{2}{p{0.3\linewidth}}{
        \cellcolor[rgb]{0.9,0.9,0.9}{
            \makecell[{{p{\linewidth}}}]{
                \texttt{\tiny{[GM$|$GM]}}
                [1, 0, 1]
            }
        }
    }
    & & \\ \\

    \theutterance \stepcounter{utterance}  
    & & \multicolumn{4}{p{0.6\linewidth}}{
        \cellcolor[rgb]{0.9,0.9,0.9}{
            \makecell[{{p{\linewidth}}}]{
                \texttt{\tiny{[GM$\rangle$P2]}}
                \texttt{Den anden spiller har indsendt et hemmeligt forslag. Nu er det din tur til at indsende dit eget forslag.} \\
            }
        }
    }
    & \\ \\

    \theutterance \stepcounter{utterance}  
    & & & \multicolumn{4}{p{0.6\linewidth}}{
        \cellcolor[rgb]{0.9,0.9,0.9}{
            \makecell[{{p{\linewidth}}}]{
                \texttt{\tiny{[GM$\langle$P2]}}
                \texttt{{[}Forslag: 3 lampe, 1 pose{]}} \\
            }
        }
    }
     \\ \\

    \theutterance \stepcounter{utterance}  
    & & & \multicolumn{2}{p{0.3\linewidth}}{
        \cellcolor[rgb]{0.9,0.9,0.9}{
            \makecell[{{p{\linewidth}}}]{
                \texttt{\tiny{[GM$|$GM]}}
                \texttt{proposal} \\
            }
        }
    }
    & & \\ \\

    \theutterance \stepcounter{utterance}  
    & & & \multicolumn{2}{p{0.3\linewidth}}{
        \cellcolor[rgb]{0.9,0.9,0.9}{
            \makecell[{{p{\linewidth}}}]{
                \texttt{\tiny{[GM$|$GM]}}
                [1, 3, 0]
            }
        }
    }
    & & \\ \\

    \theutterance \stepcounter{utterance}  
    & & & \multicolumn{2}{p{0.3\linewidth}}{
        \cellcolor[rgb]{0.9,0.9,0.9}{
            \makecell[{{p{\linewidth}}}]{
                \texttt{\tiny{[GM$|$GM]}}
                [[1, 0, 1], [1, 3, 0]]
            }
        }
    }
    & & \\ \\

\end{supertabular}
}

\end{document}
